\subsection{DEPENDENCE ON INITIAL CONDITIONS}

Let \(\mathbf{x}^{\prime}=\mathbf{f}(t,\mathbf{x})\) be a locally Lipschitz equation on \(I\times H\) ; by th. 2 (IV, p.~172), for every point \((t_{0},\mathbf{x}_{0})\) of \(I\times H\) there exists a largest interval \(J(t_{0},\mathbf{x}_{0})\subset I\) , not reducing to a single point, containing \(t_{0}\) , and on which there exists an integral (and only one) of this equation, equal to \(x_{0}\) at \(t_{0}\) ; we shall clarify the manner in which this integral, and the interval \(J(t_{0},\mathbf{x}_{0})\) where it is defined, depend on the point \((t_{0},\mathbf{x}_{0})\) .

THEOREM 4. Let \(\mathbf{f}\) be a locally Lipschitz function on \(\mathrm{I} \times \mathrm{H}\) and \((a, \mathbf{b})\) an arbitrary point of \(\mathrm{I} \times \mathrm{H}\) .

\(1^{\circ}\) There exist an interval \(\mathbf{K} \subset \mathbf{I}\) , a neighbourhood of \(a\) in \(\mathbf{I}\) , and a neighbourhood \(\mathbf{V}\) of \(\mathbf{b}\) in \(\mathbf{H}\) such that, for every point \((t_0, \mathbf{x}_0)\) of \(\mathbf{K} \times \mathbf{V}\) , there exists an integral (and only one) \(\mathbf{u}(t, t_0, \mathbf{x}_0)\) defined on \(\mathbf{K}\) , with values in \(\mathbf{H}\) and equal to \(\mathbf{x}_0\) at the point \(t_0\) (in other words, \(J(t_0, \mathbf{x}_0) \supset \mathbf{K}\) for all \((t_0, \mathbf{x}_0) \in \mathbf{K} \times \mathbf{V}\) ).

\(2^{\circ}\) The map \((t,t_0,\mathbf{x}_0)\mapsto \mathbf{u}(t,t_0,\mathbf{x}_0)\) of \(\mathrm{K}\times \mathrm{K}\times \mathrm{V}\) into H is uniformly continuous.

\(3^{\circ}\) There exists a neighbourhood \(\mathbf{W} \subset \mathbf{V}\) of \(\mathbf{b}\) in \(\mathbf{H}\) such that, for every point

\[
(t, t _ {0}, \mathbf {x} _ {0}) \in \mathrm{K} \times \mathrm{K} \times \mathrm{W},
\]

the equation \(\mathbf{x}_{0} = \mathbf{u}(t_{0}, t, \mathbf{x})\) has a unique solution x on V equal to \(\mathbf{u}(t, t_{0}, \mathbf{x}_{0})\) (``resolution of the integral with respect to the constant of integration'').

\(1^{\circ}\) Let S be a ball with centre \(\mathbf{b}\) and radius \(r\) contained in H, and \(\mathrm{J}_0\) an interval contained in I, a neighbourhood of \(a\) in I, such that \(\mathbf{f}\) is bounded and Lipschitz (with respect to some constant \(k\) ) on \(\mathrm{J}_0 \times \mathrm{S}\) ; denote by M the supremum of \(\| \mathbf{f}(t, \mathbf{x})\|\) on \(\mathrm{J}_0 \times \mathrm{S}\) . Then there exist (IV, p.~171, th. 1) an interval \(\mathrm{J} \subset \mathrm{J}_0\) , a neighbourhood of \(a\) in I, and an integral \(\mathbf{v}\) of \(\mathbf{x}' = \mathbf{f}(t, \mathbf{x})\) defined on J, with values in S and equal to \(\mathbf{b}\) at \(a\) . We shall see that the open ball V with centre \(\mathbf{b}\) and radius \(r/2\) , and the intersection K of J with an interval \([a - l, a + l[\) , where \(l\) is small enough, are as required. Indeed, prop. 7 of IV, p.~174 (applied to the set \(\mathrm{J}_0 \times \mathrm{S}\) and the case where \(\alpha = 0\) ) shows that there exists an integral of \(\mathbf{x}' = \mathbf{f}(t, \mathbf{x})\) defined on K, with values in S, and equal to \(\mathbf{x}_0\) at a point \(t_0 \in \mathrm{K}\) , provided that

\[
\left\| \mathbf {v} (t) - \mathbf {b} \right\| + \left\| \mathbf {v} (t _ {0}) - \mathbf {x} _ {0} \right\| e ^ {k | t - t _ {0} |} <   r\tag{18}
\]

for every \(t \in \mathbf{K}\) . Now, by the mean value theorem, one has

\[
\| \mathbf {v} (t) - \mathbf {b} \| \leqslant \mathrm{M} | t - a | \leqslant \mathrm{M} l
\]

for every \(t \in \mathbf{K}\) ; since \(\| \mathbf{x}_0 - \mathbf{b} \| < r / 2\) one sees that it suffices to take \(l\) such that

\[
\mathbf {M} l + (\mathbf {M} l + r / 2) e ^ {2 k l} <   r\tag{19}
\]

for the relation (18) to be satisfied for every \((t, t_{0}, \mathbf{x}_{0})\) of \(K \times K \times V\) .

\(2^{\circ}\) By the mean value theorem we have

\[
\left\| \mathbf {u} (t _ {1}, t _ {0}, \mathbf {x} _ {0}) - \mathbf {u} (t _ {2}, t _ {0}, \mathbf {x} _ {0}) \right\| \leqslant \mathrm{M} \left| t _ {2} - t _ {1} \right|\tag{20}
\]

for all \(t_0, t_1, t_2\) in K and \(\mathbf{x}_0\) in V. Now prop. 5 (IV, p.~169) shows that

\[
\left\| \mathbf {u} (t, t _ {0}, \mathbf {x} _ {1}) - \mathbf {u} (t, t _ {0}, \mathbf {x} _ {2}) \right\| \leqslant e ^ {2 k t} \left\| \mathbf {x} _ {2} - \mathbf {x} _ {1} \right\|\tag{21}
\]

for every t and \(t_{0}\) in K, and \(x_{1}\) and \(x_{2}\) in V. Finally, if \(t_{1}\) and \(t_{2}\) are any two points in K,

\[
\left\| \mathbf {u} \left(t _ {1}, t _ {2}, \mathbf {x} _ {0}\right) - \mathbf {u} \left(t _ {2}, t _ {2}, \mathbf {x} _ {0}\right) \right\| \leqslant \mathrm{M} \left| t _ {2} - t _ {1} \right|
\]

by the mean value theorem, that is to say

\[
\left\| \mathbf {u} (t _ {1}, t _ {2}, \mathbf {x} _ {0}) - \mathbf {x} _ {0} \right\| \leqslant \mathbf {M} \left| t _ {2} - t _ {1} \right|;
\]

since \(\mathbf{u}(t, t_2, \mathbf{x}_0)\) is identical to the integral which takes the value \(\mathbf{u}(t_1, t_2, \mathbf{x}_0)\) at the point \(t_1\) , prop. 5 (IV, p.~169) shows that

\[
\left\| \mathbf {u} (t, t _ {1}, \mathbf {x} _ {0}) - \mathbf {u} (t, t _ {2}, \mathbf {x} _ {0}) \right\| \leqslant \mathrm{Me} ^ {2 k l} | t _ {2} - t _ {1} |\tag{22}
\]

for all \(t, t_{1}, t_{2}\) in \(\mathbf{K}\) and \(\mathbf{x}_{0}\) in \(\mathrm{V}\) . The three inequalities (20), (21) and (22) thus prove the uniform continuity of the map \((t, t_{0}, \mathbf{x}_{0}) \mapsto \mathbf{u}(t, t_{0}, \mathbf{x}_{0})\) on \(\mathrm{K} \times \mathrm{K} \times \mathrm{V}\) .

\(3^{\circ}\) By (20), we have \(\|\mathbf{u}(t, t_{0}, \mathbf{x}_{0}) - \mathbf{x}_{0}\| \leqslant \mathbf{M} |t - t_{0}| \leqslant 2\mathbf{M}l\) on

\[
\mathbf {K} \times \mathbf {K} \times \mathbf {V}.
\]

If \(l\) is taken small enough, so that \(2\mathrm{M}l < r / 4\) , one then sees that if \(\mathbf{x}_0\) is any point of the open ball \(\mathbf{W}\) with centre \(\mathbf{b}\) and radius \(r / 4\) , that \(\mathbf{u}(t,t_0,\mathbf{x}_0)\in \mathbf{V}\) for any \(t\) and \(t_0\) in \(\mathbf{K}\) . If \(\mathbf{x} = \mathbf{u}(t,t_0,\mathbf{x}_0)\) , the function \(s\mapsto \mathbf{u}(s,t,\mathbf{x})\) is then defined on \(\mathbf{K}\) and is equal to the integral of (1) which takes the value \(\mathbf{x}\) at the point \(t\) , that is, to \(\mathbf{u}(s,t_0,\mathbf{x}_0)\) ; in particular

\[
\mathbf {x} _ {0} = \mathbf {u} (t _ {0}, t _ {0}, \mathbf {x} _ {0}) = \mathbf {u} (t _ {0}, t, \mathbf {x}).
\]

Moreover, if \(y \in V\) is such that \(\mathbf{x}_{0} = \mathbf{u}(t_{0}, t, \mathbf{y})\) , then the integral \(s \mapsto \mathbf{u}(s, t, \mathbf{y})\) takes the value \(x_{0}\) at \(t_{0}\) so is identical to \(s \mapsto \mathbf{u}(s, t_{0}, \mathbf{x}_{0})\) , which consequently takes the value x at t, which shows that y = x and concludes the proof.

